%%%%%%%%%%%%%%%%%%%%%%%%%%%%%%%%%%%%%%%%%%%%%%%%%
%%%%%%%%%%%% chap: conclusions %%%%%%%%%%%%%%%%%%%%
%%%%%%%%%%%%%%%%%%%%%%%%%%%%%%%%%%%%%%%%%%%%%%%%%

\chapter{Conclusions}\label{chapter:chap5}

\section{Summary}

This thesis presented the design, implementation, and evaluation of a
multi-stage phishing honeypot capable of capturing the complete attacker
lifecycle — from initial credential theft through to post-compromise SSH
activity — and classifying attacker behaviour using machine learning.

The system integrates six components: a MailHog-based email server, a
three-stage Flask phishing decoy, a Cowrie SSH honeypot, a PostgreSQL
database, an ELK log pipeline, and a simulation layer comprising four
bot profiles. A \texttt{session\_id} propagated across all components
links each phishing visit to its downstream SSH session, enabling
multi-stage behavioural analysis that is absent from single-stage
honeypot deployments.

Four bot profiles — scanner, credential stuffer, human-like, and
sophisticated — were implemented with timing distributions derived from
the CyberLab real-world honeynet dataset \cite{Alata2006Honeynet} and
command sequences grounded in MITRE ATT\&CK techniques T1078, T1059,
T1087, and T1005. Each profile was assigned a dedicated Docker network
subnet, giving it a unique source IP address, allowing ground-truth
labels to be set with certainty.

Two Random Forest classifiers were trained on 3,141 collected sessions.
RF1 classifies attacker profile (\textit{bot type}) across four classes
and achieves 97.6\% accuracy and 98.6\% macro F1 in 5-fold stratified
cross-validation. RF2 classifies attack stage (credential submission
versus SSH attack) and achieves perfect accuracy and F1, confirming
that the structural encoding of multi-stage sessions is a reliable
discriminating signal. These results meet or exceed published Random
Forest performance on the CICIDS2017 benchmark
\cite{Bringer2012HoneypotSurvey}, despite using only six
application-layer features per session rather than network flow
statistics.

\section{Contributions}

The principal contributions of this work are:

\begin{enumerate}
    \item \textbf{Multi-stage session linkage.} The use of a single
    \texttt{session\_id} to bridge phishing and SSH events produces a
    richer behavioural record than either stage alone. This design
    choice enables classifiers to exploit the cross-stage fingerprint
    and is, to the best of the author's knowledge, not present in
    prior academic honeypot datasets.

    \item \textbf{Realistic synthetic traffic generation.} The bot
    simulation reproduces empirically-grounded timing behaviour, deriving
    delays from a 293-session real-world dataset rather than from
    arbitrary parameters. This justification makes the synthetic dataset
    a plausible proxy for real attacker populations.

    \item \textbf{Attacker profiling from behavioural features.} The
    demonstrated sufficiency of six timing and interaction features for
    high-accuracy four-class profiling suggests that application-layer
    honeypot logs, which are cheaper to collect than full packet captures,
    are a viable basis for automated threat classification.

    \item \textbf{Integrated detection pipeline.} The end-to-end system
    — from email decoy to Elasticsearch indexing and Kibana
    visualisation — provides a deployable honeypot architecture that
    can be extended to real-world attacker traffic with minimal
    modification.
\end{enumerate}

\section{Limitations}

Several limitations constrain the generalisability of the results.

First, the dataset is entirely synthetic. Although the bot timing
distributions are derived from real-world data, the bots execute
deterministic profiles that produce cleanly separable feature clusters.
Real attackers are less predictable and more diverse, and a classifier
trained on this data may not transfer directly to a production
environment without retraining on live traffic.

Second, the classifier relies on six features, two of which
(\texttt{field\_corrections} and \texttt{inter\_cmd\_delay\_avg\_s})
may be spoofed by a sophisticated adversary who is aware of the
honeypot. A real attacker who deliberately varies their inter-command
delay or introduces artificial corrections could evade classification.

Third, the CICIDS2017 comparison is illustrative rather than direct.
The two datasets measure different properties — application-layer
interaction timing versus network flow statistics — and address
different classification tasks. The comparison contextualises the
results but does not constitute a controlled experimental evaluation
against a shared test set.

\section{Future Work}

Several directions extend the current work.

\textbf{Live deployment and retraining.} Deploying the honeypot on a
publicly accessible network would accumulate real attacker sessions,
enabling the classifier to be retrained on genuine data. Periodic
retraining with active learning could maintain classification accuracy
as attacker behaviour evolves.

\textbf{Additional attacker profiles.} The current four profiles cover
a spectrum from fully automated to deliberate human-speed interaction.
Extending the simulation to include exploit-kit behaviour, credential
validation loops, and lateral movement stages would increase the
realism and forensic utility of the dataset.

\textbf{Natural language processing of command sequences.} The current
model treats command activity only through aggregate statistics
(\texttt{command\_count}, \texttt{inter\_cmd\_delay\_avg\_s}). Applying
sequence-level models — such as recurrent networks or transformer-based
embeddings over the command token sequence — could extract richer
behavioural signals, particularly for distinguishing the human-like and
sophisticated profiles that share similar timing ranges.

\textbf{Adversarial robustness.} Evaluating classifier robustness
against adversarially perturbed inputs — for example, a scanner that
artificially inflates its page dwell time — would quantify the attack
surface of the current feature set and motivate hardened features.

\textbf{Multi-honeypot federation.} Aggregating sessions across
multiple geographically distributed honeypot instances would capture
attacker population diversity and enable longitudinal analysis of
campaign-level behaviour patterns.

\section{Closing Remarks}

The results demonstrate that a multi-stage phishing honeypot capturing
both application-layer and SSH behaviour produces a dataset that is
sufficient for accurate automated attacker profiling. The core insight
— that timing and interaction behaviour constitute a reliable
cross-stage fingerprint — has practical implications for threat
intelligence: organisations that instrument credential-harvesting pages
alongside their SSH infrastructure can extract actionable attacker
classification from the resulting logs without requiring network-level
capture.

The system architecture and dataset generation methodology are
reproducible from the source code and configuration provided with this
thesis, and are intended to serve as a baseline for future work on
multi-stage honeypot analysis.
